% 17-termination.tex — Chapter 17: Termination Verification
\chapter{Termination Verification}
\label{chap:17-termination}
\index{termination verification}
\index{bounded while}
\index{decreasing measure}

\pnum
Lain's \lain{func} declarations are required to terminate.  Termination is
guaranteed by the combination of three rules:
\begin{enumerate}
\item Direct recursion is rejected in \lain{func} at the resolve phase.
\item Mutual recursion among \lain{func} declarations is rejected by
  \cname{sema\_check\_no\_mutual\_recursion} (\cref{sec:sema-mrec}).
\item Unbounded \lain{while} loops are forbidden in \lain{func}.
  A \lain{while} loop with a \lain{decreasing measure} is allowed because the
  compiler verifies termination.
\end{enumerate}

\pnum
The bounded-while verifier is self-contained in \srcref{src/sema.h:344} and
does \emph{not} use the VRA range table, because VRA cannot track measures that
involve struct fields.

% -------------------------------------------------------------------------
\section{Termination measure requirements}
\label{sec:term-measure}
\index{termination measure}
% -------------------------------------------------------------------------

\pnum
A valid termination measure must satisfy two conditions:
\begin{enumerate}
\item \textbf{Non-negativity}: when the loop condition holds, the measure is $\geq 0$.
\item \textbf{Strict decrease}: the loop body strictly decreases the measure.
\end{enumerate}

\pnum
The verifier checks these conditions syntactically and raises errors if they
cannot be proven:

\begin{center}
\begin{tabular}{lll}
\toprule
\textbf{Error} & \textbf{Condition} & \textbf{Hint} \\
\midrule
\diagcode{E080} & Non-negativity cannot be proven & Use \lain{while a < b decreasing b - a} \\
\diagcode{E081} & No variables in measure & Measure must reference identifiers/fields \\
\diagcode{E082} & No strict decrease found & Body must contain an assignment that decreases the measure \\
\bottomrule
\end{tabular}
\end{center}

% -------------------------------------------------------------------------
\section{Non-negativity verification}
\label{sec:term-nonneg}
% -------------------------------------------------------------------------

\pnum
\cname{sema\_verify\_measure\_nonneg(cond, measure)} checks whether the loop
condition syntactically implies the measure is non-negative.  It recognizes
two patterns:

\begin{enumerate}
\item \textbf{Difference pattern}: measure has the form \lain{hi - lo} and the
  condition contains a comparison \lain{lo < hi} (strict) or \lain{lo <= hi}.
  Example: \lain{while a < b decreasing b - a}.
\item \textbf{Single-variable pattern}: measure is a single identifier \texttt{n}
  and the condition contains \lain{n > 0} or \lain{n > constant>=0}.
  Example: \lain{while n > 0 decreasing n}.
\item \textbf{In-guard pattern}: condition is \lain{idx in arr} and measure
  is \lain{arr.len - idx}.  The \lain{in} guard implies
  $0 \leq \texttt{idx} < \texttt{arr.len}$, so the measure is $\geq 1 > 0$.
\end{enumerate}

\pnum
Conjunctive conditions (\lain{a and b}) are traversed recursively; either
conjunct may supply the non-negativity proof.

% -------------------------------------------------------------------------
\section{Strict decrease verification}
\label{sec:term-decrease}
% -------------------------------------------------------------------------

\pnum
\cname{measure\_scan\_body(body, vars, nvar)} scans the loop body for assignments
to variables that appear in the measure expression.  \cname{measure\_extract\_vars}
first decomposes the measure into (variable, polarity) pairs:
\begin{itemize}
\item Measure \texttt{x}: variable \texttt{x} with polarity $+1$.
\item Measure \texttt{a - b}: variable \texttt{a} with polarity $+1$,
  variable \texttt{b} with polarity $-1$.
\end{itemize}

\pnum
For each assignment \lain{v = v op k} (where \texttt{v} is a measure variable
and \texttt{k} is a literal), \cname{assignment\_direction} determines whether
the assignment increases or decreases \texttt{v}.  The effect on the measure is
\texttt{polarity * direction}:
\begin{itemize}
\item Negative effect (direction opposes polarity): the measure decreases; record
  \texttt{found\_decrease = true}.
\item Positive effect (direction reinforces polarity): the measure \emph{increases};
  immediately return conflict ($-1$).
\end{itemize}

\pnum
If no assignment to any measure variable is found (\texttt{found\_decrease} remains
false), the verifier returns 0 and error \diagcode{E082} is emitted.

% -------------------------------------------------------------------------
\section{Accepted measure patterns}
\label{sec:term-patterns}
% -------------------------------------------------------------------------

\pnum
Practical examples of accepted termination measures:

\begin{laincode}
// while n > 0 decreasing n { n -= 1 }
while n > 0 decreasing n {
    n -= 1
}

// while pos < size decreasing size - pos { pos += 1 }
while pos < size decreasing size - pos {
    pos += 1
}

// while idx in arr decreasing arr.len - idx { idx += 1 }
while idx in arr decreasing arr.len - idx {
    idx += 1
}
\end{laincode}

\pnum
All three patterns are accepted because: (1) the condition implies the measure is
positive, and (2) the body contains an assignment that decreases the measure by a
positive literal amount.
